\documentclass[acmsmall,screen,nonacm,11pt]{acmart}
\geometry{letterpaper,margin=1in,headheight=14pt,heightrounded=false}
\settopmatter{printacmref=false,printccs=false,printfolios=true}
\setcopyright{none}
\renewcommand{\bibfont}{\normalfont\normalsize}
\renewcommand{\shortauthors}{Zijian Zeng}
\newcommand{\eps}{\varepsilon}
\newcommand{\co}{\mathrel{<_{\mathrm{co}}}}
\newcommand{\cole}{\mathrel{\leq_{\mathrm{co}}}}
\newcommand{\rev}{\operatorname{rev}}
\newcommand{\lcp}{\operatorname{lcp}}
\newcommand{\dom}{\operatorname{dom}}
\newcommand{\Pref}{\operatorname{Pref}}
\newcommand{\Rep}{\operatorname{Rep}}
\newcommand{\out}{\operatorname{out}}
\newcommand{\sgnset}{\{+,-\}}
\newtheorem{theorem}{Theorem}[section]
\newtheorem{lemma}[theorem]{Lemma}
\newtheorem{corollary}[theorem]{Corollary}
\theoremstyle{definition}
\newtheorem{definition}[theorem]{Definition}
\title{Deciding Membership and Realizability for Wheeler Transducers}
\author{Zijian Zeng}
\affiliation{\institution{UCSI University}\city{Kuala Lumpur}\country{Malaysia}}
\email{1002266693@ucsiuniversity.edu.my}
\begin{abstract}
\normalsize
Wheeler transducers combine a total colexicographic order on reaching
languages with a finite convex piecewise monotonicity condition on accumulated
output. The number of monotone pieces is unspecified. We give a deterministic
polynomial-time membership algorithm for finite trimmed deterministic
subsequential transducers, allowing partial transitions and erasing outputs.
This resolves the membership question for the piecewise definition of
D'Agostino and Paradiso.

The structural criterion distinguishes two forms of recurrence. In a branching
strongly connected component, a strict output turn can be replicated in
disjoint input intervals, so finite piecewise monotonicity requires global
monotonicity. A simple cycle instead has two sides approaching a central cut;
each side is monotone exactly when two adjacent levels are monotone. These
conditions reduce membership to finitely many comparisons of independent
regular output ranges. A direct implementation takes $O(S^{13})$ bit
operations for explicit encoding size $S$; this bound is conservative.
Every output failure after input-order success has a common-loop pumped
certificate of disjoint strict turns, with a total computable template bound.

As a consequence, equivalent Wheeler realizability is decidable and a
realization can be constructed on positive instances. This uses classical
prefix normalization and residual merging, an effective automatic convex
quotient, and the membership algorithm. No complexity or minimality claim is
made for the realization construction.
\end{abstract}
\begin{document}
\maketitle
\pagestyle{plain}
\newpage

\section{The decision problems and the contribution}

A subsequential transducer reads its input deterministically, emits a finite
word on each transition, and may append a terminal word. D'Agostino and
Paradiso define a Wheeler transducer by requiring both a total colexicographic
order on its reaching languages and finite piecewise monotonicity of its
accumulated output on each such language~\cite{DP26,DP26short}.
Their current definition permits a different finite number of convex pieces
at each state, with either orientation on each piece. Section~4 of both
versions asks whether membership is decidable and whether an equivalent
Wheeler realization can be found. The issue is effective control of an
unbounded partition, rather than recognition for a prescribed number of pieces.

Our principal contribution is an exact criterion on strongly connected
components (SCCs). After checking input order, global output monotonicity is
decidable by propagating local violations. Branching cyclic SCCs must satisfy
this global condition. A simple cycle permits a different phenomenon: two
monotone sides can meet at an unrealized cut, even with infinitely many turns
crossing that cut. We characterize its two sides by adjacent-level tests.
Acyclic states then inherit finite partitions. All tests use finite automata
for independent output ranges; no arbitrary intersection of rational
relations is required. This supplies an algorithm, a polynomial bound, and
effective obstructions for the original unrestricted membership problem.

For context, Wheeler \emph{language} recognition from a DFA already has an
$O(mn)$ algorithm~\cite{Becker23}. That problem does not test accumulated
output. Prefix normalization and minimization are classical~\cite{BC00},
as are effective first-order queries and representatives in automatic
structures~\cite{Blum99,BG00}. We use these as established tools for the
realizability consequence. The statewise residual-prefix identity used there
is already recorded in Remark~3 of~\cite{DP26}. The independent structural
step here is the SCC criterion for finite convex output partitions.

\subsection{Model and statements}

Both finite alphabets $\Sigma,\Gamma$ have specified total orders. For words,
$u\co v$ means $\rev(u)<_{\mathrm{lex}}\rev(v)$, with a proper prefix smaller
in lexicographic order. Thus a proper suffix is smaller in colexicographic
order. A set is convex relative to an ordered domain if it contains every
domain element between two of its elements. A function is \emph{piecewise
monotone} if its domain has a finite partition into nonempty convex sets on
each of which it is weakly increasing or weakly decreasing.

Let $T=(Q,s,\Sigma,\Gamma,\delta,\omega,F,t)$ be finite and trimmed:
every state is accessible and can reach a final state. The deterministic
transition function $\delta$ is partial; $\omega$ has the same domain and
takes values in $\Gamma^*$, including $\eps$. There is no initial output.
For a defined path, write $A(u)=\omega(s,u)$ and
\[
 I_q=\{u:\delta(s,u)=q\},\qquad D=\bigcup_{q\in Q}I_q.
\]
On accepted inputs the partial function is $f(u)=A(u)t(\delta(s,u))$.
Trimming gives $D=\Pref(\dom f)$. A nonempty trimmed instance has a
nonempty accepted domain. An encoding admitting an empty function can handle
that degenerate case separately.

We use precisely the literal reaching-language definition in
Definitions~4--5 of~\cite{DP26}: distinct states must satisfy either
$I_p<I_q$ or $I_q<I_p$, where $X<Y$ abbreviates
$\forall x\in X\,\forall y\in Y\;(x\co y)$; and $A|_{I_q}$ must be
piecewise monotone for every $q$. No additional input-consistency or
initial-indegree convention is imposed. Membership tests accumulated
transition output $A$, not the terminal-output function $f$.

\begin{theorem}\label{thm:member}
Membership in this class is decidable in $O(S^{13})$ bit operations, where
$S\geq2$ is the total explicit encoding size, counting states, transitions,
ordered alphabet descriptions, and all output characters. Partial
transitions, erasure, and arbitrary terminal words are allowed.
\end{theorem}

\begin{theorem}\label{thm:realize}
Given such a transducer, it is decidable whether an equivalent finite Wheeler
transducer over the same ordered alphabets exists. On a positive instance one
can construct one, without introducing an initial output. No polynomial bound
or minimality assertion is made for this construction.
\end{theorem}

Section~\ref{sec:cert} proves the additional pumped-certificate assertion.
All arguments below apply to infinite reaching languages. Finite controls
described at the end supplement the proofs rather than establishing them.

\section{Convex blocks and exact global monotonicity}\label{sec:blocks}

We first check whether all reaching languages are pairwise separated. They
are regular. Reverse their automata and use the regular-language comparison
test below to decide each of the two possible universal strict orders.
Reject a pair for which neither holds. Otherwise the induced state relation
is a strict total order: transitivity follows directly from the universal
definition and nonemptiness. Each $I_q$ is then convex in $D$, since an
intervening word in another state would contradict separation.

\begin{lemma}[Recursive blocks]\label{lem:blocks}
Assume input order passes. The immediate ordered disjoint blocks of $I_q$
are $\{\eps\}$ if $q=s$, followed by all $I_p a$ with $\delta(p,a)=q$,
sorted first by $a$ and then by the order of $p$. Each is convex in $I_q$.
More generally, if $\delta(p,z)=q$, then $I_pz$ is convex in $I_q$, and
\begin{equation}\label{eq:transport}
 A(uz)=A(u)\omega(p,z)\quad(u\in I_p).
\end{equation}
Common right concatenation preserves and reflects every colex comparison.
\end{lemma}
\begin{proof}
Every nonempty word has a unique last transition. Different last letters
determine their colex order, and a common last letter cancels, leaving the
separated predecessor languages. The empty word is least. This proves the
immediate decomposition and convexity. Iterating it proves the assertion for
$z$; the output identity follows from the path definition. Finally, reversal
turns a common right suffix into a common left prefix, which cancels in lex
order. This also holds for an empty suffix and proper-suffix comparisons.
\end{proof}

Let $O_p=\{A(u):u\in I_p\}$. This is an effective regular language: start
at $s$, accept at $p$, forget input labels, and read each transition's output
word, subdividing word edges and retaining epsilon edges. A block $I_pz$
has output range $O_p\omega(p,z)$. Singleton blocks have singleton ranges.
These ranges can be used independently when the input blocks are distinct
and already ordered; no relation between their chosen words is required.

\begin{lemma}[Regular range comparison]\label{lem:range}
For two epsilon-NFAs one can decide whether some accepted words $x,y$
satisfy $x<_{\mathrm{lex}}y$, $x>_{\mathrm{lex}}y$, or $x=y$, in polynomial
time. The same holds in colex order and for universal weak or strict order.
\end{lemma}
\begin{proof}
Compute epsilon closures. Synchronize both automata on equal symbols until
the first mismatch. A strict inequality is possible exactly when a reachable
product state either permits two suitably ordered distinct next symbols
with accepting continuations on both sides, or permits acceptance on the
shorter side and a positive-length accepted continuation on the other.
All these predicates are finite reachability tests. Equality uses simultaneous
acceptance after a common word. Reversing automata gives colex comparison.
Universal weak order excludes the opposite strict comparison; universal
strict order excludes both the opposite comparison and equality.
\end{proof}

For $e\in\sgnset$, orientation $+$ means weakly increasing and $-$ weakly
decreasing. Let $S_e$ be the states with a violating pair from two distinct
immediate blocks in their prescribed order. Lemma~\ref{lem:range} decides
these seeds. Let $B_e$ contain all states reachable in the input graph from
$S_e$, including seeds, and put $G_e(q)\iff q\notin B_e$.

\begin{lemma}\label{lem:global}
$G_e(q)$ holds exactly when $A|_{I_q}$ has global orientation $e$.
\end{lemma}
\begin{proof}
A seed violation propagates along every defined path by
\eqref{eq:transport}. Conversely, take a violating pair at $q$. If its words
are in different immediate blocks, $q$ is a seed. Otherwise they lie in one
$I_pa$; neither can be the singleton epsilon block. Delete the common last
letter. The pair remains strictly violating at $p$. Repetition terminates
because the sum of lengths decreases, reaching a seed. The removed letters
give a path from that seed to $q$.
\end{proof}

Write $M(q)=G_+(q)\lor G_-(q)$. We will also use two elementary order facts.
First, a map between total orders is globally monotone in some orientation
exactly when it has no strict peak or valley triple. Indeed, if a rising and
a falling pair exist, sort their at most four arguments. Delete consecutive
equal values. The successive strict signs include both orientations, so a
sign change gives a strict turn. Second, infinitely many strictly separated
closed strict-turn intervals preclude every finite convex monotone
partition: each interval must cross a piece boundary, and a boundary can
split at most one interval in such a family. Both increasing and decreasing
families of intervals are covered.

\section{The SCC criterion}\label{sec:scc}

SCCs refer to the ordinary input graph, counting distinct labeled edges.
A cyclic SCC is \emph{simple} if its internal graph is a directed simple
cycle, including a single self-loop. A loop-free singleton is acyclic.

\subsection{Branching recurrence}

\begin{lemma}\label{lem:branch}
In a cyclic nonsimple SCC, finite piecewise monotonicity at every state is
equivalent to $M(q)$ at every state.
\end{lemma}
\begin{proof}
Sufficiency is immediate. For necessity, a nonsimple strongly connected
component has a vertex $r$ with two internal outgoing edges. Determinism
makes their input letters distinct. Completing return paths gives nonempty
loops $u,v$ at $r$ with distinct first letters, hence $uv\ne vu$.
The infinite words $\rev(u)^\omega$ and $\rev(v)^\omega$ must differ:
if equal to $w$, the two periods imply that both
$\rev(u)\rev(v)$ and $\rev(v)\rev(u)$ are the same-length prefix of $w$.
Effectively choose $k$ large enough that $A_0=u^k$ and $B_0=v^k$ have a
mismatch in reversal before either ends. They are suffix-incomparable.

If some $q$ in the component fails $M(q)$, it has a strict turn. Appending
a path from $q$ to $r$ transports that turn to words
$\alpha_1\co\alpha_2\co\alpha_3$ reaching $r$. For $n\geq0$ set
\begin{equation}\label{eq:branchpump}
 U_i(n)=\alpha_i A_0B_0^n\quad(1\leq i\leq3).
\end{equation}
The outputs retain the same strict turn by common suffix cancellation.
The suffix cylinders of $A_0B_0^n$ are strictly separated in a fixed direction.
For consecutive indices, cancel the common reversed $B_0^n$; the comparison
then encounters the fixed mismatch between $\rev(A_0)$ and $\rev(B_0)$,
before any prefix $\alpha_i$ can matter. Thus either every interval precedes
the next, or every interval follows it. The order fact above rules out a
finite partition at $r$, proving necessity for the component property.
\end{proof}

\subsection{Simple cycles and their two sides}

Let $C$ be a simple cyclic SCC. Choose $q\in C$, let $c$ be the nonempty
input word for one full traversal from $q$ to itself, and put
$b=\omega(q,c)$. Expand the unique internal predecessor block repeatedly
for $|C|$ steps in Lemma~\ref{lem:blocks}. Retain that child and list all
other children to its left or right. This yields
\begin{equation}\label{eq:shell}
 I_q=L\;\dot\cup\;I_qc\;\dot\cup\;R,\qquad L<I_qc<R.
\end{equation}
Each side is a finite ordered list of singleton words arising from $s$, or
blocks $I_pz$ with $p\notin C$ and $z$ an explicit entering path of length
at most $|C|$. There is only one internal predecessor at every step; this
is why the construction stops after one traversal. Either side may be empty.

Put $L_n=Lc^n$ and $R_n=Rc^n$. Iteration of~\eqref{eq:shell} gives
\begin{equation}\label{eq:levels}
 I_q=\Bigl(\dot\bigcup_{n\geq0}L_n\Bigr)
       \;\dot\cup\;
       \Bigl(\dot\bigcup_{n\geq0}R_n\Bigr).
\end{equation}
There is no leftover finite word: membership in every $I_qc^n$ would require
length at least $n|c|$ for all $n$. The $L_n$ increase in the order, the
$R_n$ decrease, every left level precedes every right level, and each level
and each whole side is convex in $I_q$.

\begin{lemma}[Two-side test]\label{lem:cycle}
$A|_{I_q}$ is piecewise monotone if and only if each restriction
\begin{equation}\label{eq:adjacent}
 A|_{L\cup Lc}\quad\hbox{and}\quad A|_{Rc\cup R}
\end{equation}
is globally weakly monotone, with directions chosen independently.
If this holds, two pieces suffice at $q$.
\end{lemma}
\begin{proof}
If one restriction fails both orientations, it has a strict turn
$v_1\co v_2\co v_3$. Appending $c^{2n}$ places its three points on levels
$2n$ and/or $2n+1$ of the same side. Consecutive pairs of levels are strictly
separated, increasingly on the left and decreasingly on the right.
The strict output turn persists after appending $b^{2n}$, including when
$b=\eps$. Infinitely many separated turn intervals contradict any finite
partition.

Conversely, if the left restriction has orientation $e$, appending $c^n$
proves the same orientation on $L_n\cup L_{n+1}$ for every $n$. If $L$ is
nonempty, all its levels are nonempty. For points on nonadjacent levels,
insert one point in each intervening level and apply transitivity of the
weak output order. Hence the whole left side has orientation $e$. An empty
side needs no piece. The identical argument in decreasing level order
applies on the right, using its independently chosen orientation. The two
convex sides in~\eqref{eq:levels} give the asserted partition.
\end{proof}

The test is effective. Expand~\eqref{eq:adjacent} into its finite ordered
lists, appending $c$ to paths for the second copies. Orientation $e$ holds
on such a list exactly when every nonsingleton block $I_pz$ has $G_e(p)$,
and every earlier/later pair of blocks has output ranges in weak order $e$.
Lemma~\ref{lem:blocks} checks the within-block condition exactly; the
independent ranges and Lemma~\ref{lem:range} check all cross-block pairs.
Thus arbitrary regular prefix-output ranges are retained, without a depth
bound or a bounded-expression assumption. All states of $C$ inherit finite
partitions from $q$: append a path from the chosen state to $q$, restrict
the partition there, and use suffix cancellation. The preimages are convex.

The separation of the sides is essential. Consider the three-state machine
with transitions $s\xrightarrow{0/\eps}p$,
$s\xrightarrow{1/\eps}r$, $r\xrightarrow{0/\eps}p$,
$p\xrightarrow{0/0}p$, sole final state $p$, and empty terminal output.
Here $I_p=0^+\cup10^+$; $A(0^m)=A(10^m)=0^{m-1}$.
Its first part increases and its second decreases in colex. Infinitely many
peak triples $0^{n+1},0^{n+2},10^{n+1}$, for $n\geq0$, cross the same
central cut, but the map has two pieces. A cyclic turn alone is therefore
not a nonmembership certificate. Lemma~\ref{lem:cycle} detects precisely
the turns that can recur in disjoint levels.

\subsection{Algorithm and proof of Theorem~\ref{thm:member}}

Check input order, compute regular ranges and the two global flags, and
compute SCCs. Require $M(q)$ at every state in each branching cyclic SCC.
At one representative of each simple cyclic SCC require the two-side test.
Reject on a failed requirement, and otherwise accept. Every rejection is
sound by the preceding lemmas. Every passing cyclic SCC has finite partitions
at all its states. The remaining components are loop-free singletons.
In a topological order, their immediate blocks inherit the finite partitions
of predecessor states, transported by their last transitions. Concatenating
these partitions, with the initial singleton if present, gives a finite
convex monotone partition. This proves sufficiency at every state.
All constructions terminate; terminal words are unused. Empty input or
output alphabets and the epsilon-only domain obey the same finite tests.

For completeness, a loose explicit resource bound follows. Write
$n=|Q|\leq S$, $m\leq S$ for the edge count, and $M\leq S$ for the maximum
edge-output length. The range NFAs have $O(S)$ states and edges. There are
at most $m+1$ immediate blocks over all states. In one annular expansion per
simple SCC, each external incoming edge occurs once, and the initial
singleton occurs at most once overall. The doubled annular lists therefore
have $O(S)$ blocks in total. Appended paths have length at most $2n$ and
output length at most $2nM$. Every compared output NFA has size $O(S^2)$,
and there are $O(S^2)$ comparisons; input-order tests use $O(n^2)$ pairs of
$O(S)$-size input automata.

With $K$ states and edges, straightforward epsilon elimination costs at
most $O((|\Gamma|+1)K^4)$ operations. Enumerating synchronized transitions,
live unequal-symbol transitions, and prefix cases in a product comparison
costs at most $O((|\Gamma|+1)^2K^4)$. For $K=O(S^2)$ this is $O(S^{10})$
per comparison, hence $O(S^{12})$ operations overall. Allowing $S$ bit
operations per explicitly encoded symbol or index comparison gives
$O(S^{13})$ bit operations. This deliberately conservative bound concerns
membership only and proves Theorem~\ref{thm:member}.

\section{Effective pumped failure certificates}\label{sec:cert}

\begin{corollary}\label{cor:pump}
For every input-ordered nonmember, some state has words
$U_i(n)=x_i y_i^n z_i$, $1\leq i\leq3$, such that for every $n\geq0$:
all three reach that state, $U_1(n)\co U_2(n)\co U_3(n)$, their outputs
form a strict peak for every $n$ or a strict valley for every $n$, and either
\[
 U_3(n)\co U_1(n+1)\ \text{for all }n,
 \quad\text{or}\quad U_3(n+1)\co U_1(n)\ \text{for all }n.
\]
The three pumps can be chosen equal. There is a total computable bound
$B(T)$ on the lengths of the nine template words needed for a certificate.
\end{corollary}
\begin{proof}
Every output rejection occurs in a branching or simple cyclic SCC.
Equation~\eqref{eq:branchpump} gives $x_i=\alpha_iA_0$, common $y_i=B_0$,
and $z_i=\eps$. The failed side in Lemma~\ref{lem:cycle} gives $x_i=v_i$,
common $y_i=c^2$, and $z_i=\eps$. The proofs supply all universal
inequalities, in either interval direction. Conversely any such certificate
precludes a finite partition by separated turn intervals.

Run the membership algorithm first. On a rejection after input order,
enumerate a finite strict-turn witness in the indicated state or regular
two-level restriction; existence was proved, so this terminates. In the
branching case the finite-power mismatch search also terminates. Output
the maximum template length as $B(T)$. Set $B(T)=0$ on acceptance,
input-order failure, or a malformed or nontrimmed description. These cases
are decidable, so $B$ is total. No numerical complexity bound for $B$ or
the witness enumeration follows from the membership runtime.
\end{proof}

Validity of a fixed template is itself decidable, independently of that
construction. For each fixed $x_i,y_i,z_i$, finite orbit decomposition of
the state action of $y_i$ gives a transient and a common period. On each
residue, defined runs and their outputs have the form $P Q^t R$ with fixed
words, allowing $Q=\eps$. A position has a given character precisely when
it lies in a fixed prefix or suffix, or in the repeated block with its
offset in a fixed residue modulo $|Q|$. All lengths are affine in $t$;
when $Q=\eps$ omit that block rather than taking a zero modulus.
Strict lex comparison is a first differing position with equality before
it, or a proper-prefix comparison. Thus input/output comparisons, with
reversal for colex, are Presburger formulas using only constant
coefficients and moduli. Finite transient checks and Presburger
decidability verify every $n$, including consecutive-index separation.
There is no multiplication of variables and no claim that a Presburger
solver was executed. The polynomial membership algorithm does not use it.

\section{Deciding equivalent realizability}\label{sec:realize}

We prove Theorem~\ref{thm:realize}. The input here need not satisfy input
order. Its partial function $f$ and prefix domain $D$ are fixed. The
construction must merge redundant residuals \emph{before} taking convex
components: otherwise a redundant unary parity coloring of a constant
function would create a false obstruction.

\subsection{Normalization and canonical colors}

Let $O_p(v)$ be the accepted continuation output from $p$, including its
terminal word, and let
$\lambda_p=\lcp\{O_p(v):v\in\dom O_p\}$.
Coaccessibility makes this range nonempty. Its output language is regular.
Find one accepted output $y$ and, among the finitely many prefixes of $y$,
choose the longest $z$ for which the entire range is contained in
$z\Gamma^*$. Regular inclusion decides these tests and computes $\lambda_p$.
For a transition $p\xrightarrow{a/w}r$, $\lambda_p$ is a prefix of
$w\lambda_r$, the LCP of all its accepted continuation outputs. It is
also a prefix of $t(p)$ when $p$ is final. Normalize to edge and terminal
words
\begin{equation}\label{eq:normalize}
 n(p,a)=\lambda_p^{-1}(w\lambda_r),\qquad
 n_t(p)=\lambda_p^{-1}t(p),
\end{equation}
where inverse notation means deletion of the verified prefix. Telescoping
gives normalized residual $g_p(v)=\lambda_p^{-1}O_p(v)$, with the same
partial domain and with empty range LCP. This is classical prefix
normalization~\cite{BC00}.

Finite partition refinement computes equality of the $g_p$. Specifically,
require equal final flags and normalized terminal words, equal defined
letters and normalized edge words, and equivalent successors. These
conditions suffice by induction on accepted continuation length. They are
also necessary for equal functions: coaccessibility makes defined letters
detectable from the partial domain, and the LCP of outputs on $a$-beginning
continuations is exactly $n(p,a)$, since its target residual has empty LCP.
Left cancellation then gives equal successor residuals. The greatest such
finite relation is therefore exact. Write $[p]$ for these merged classes.

Create a distinct new initial state $e$ with no incoming edges. Its terminal
word is $f(\eps)$ if defined. For each original edge
$s\xrightarrow{a/w}r$, create
$e\xrightarrow{a/(w\lambda_r)}[r]$; elsewhere use the normalized quotient.
Delete unreachable classes, obtaining $C$. It is trimmed and computes
exactly $f$ without initial output. For every nonempty $u\in D$ its
accumulated output is
\begin{equation}\label{eq:H}
 A_C(u)=A_T(u)\lambda_{\delta(s,u)}
       =H(u):=\lcp\{f(uv):uv\in\dom f\}.
\end{equation}
At epsilon it remains $\eps$, even if $H(\eps)\ne\eps$. Nonempty returns
to the old initial state use its normalized class, not $e$. Hence initial
cycles and nonempty initial LCPs are covered.

For nonempty $u$, the $C$-state is exactly its normalized continuation
function $R_u(v)=H(u)^{-1}f(uv)$, on $u^{-1}\dom f$. Color each nonempty
word by that state, and color epsilon by the unique color $e$. The colors
are regular reaching languages. For any other trimmed realization $W$ of
$f$, with reaching set $J_p$ and continuation LCP $\lambda'_p$, the
statewise identity already used in~\cite[Remark~3]{DP26} is
\begin{equation}\label{eq:otherH}
 H(u)=A_W(u)\lambda'_p\qquad(u\in J_p).
\end{equation}
It follows by factoring $A_W(u)$ from all accepted continuation outputs.
Thus $J_p\setminus\{\eps\}$ is monochromatic in the canonical coloring,
regardless of output timing or the original machine's redundant states.

\subsection{An effective convex quotient}

Define $E(u,v)$ on $D$ by equal canonical color and no word of a different
color strictly between $u$ and $v$ in either direction. Its classes are
exactly the maximal convex monochromatic components. Transitivity follows
because the outer interval lies in the union of the intervals to a common
point. Every convex monochromatic subset lies in a single component.

We detail the effective construction to distinguish synchronous regular
relations from arbitrary rational relations. Encode words by reversal and
use right-padded convolutions. Regular color predicates, equality, lex
order, and shortlex order are synchronous regular. A lex comparator
remembers the first differing symbol, treating an ended word as smaller;
shortlex compares lengths first and then lex on equal lengths. Ended bits
forbid a genuine symbol after padding. Boolean operations are relative
to the regular universe of valid padded tuples. Existential projection
guesses the removed track, including its end. Its witness can be longer
than every free word: after projection, take the right quotient by an
arbitrary all-padding suffix, equivalently marking a state final if it
can reach acceptance along all-padding columns. This finite construction
preserves padding invariance and imposes no length bound.

Apply these operations to
\[
 z\in D,\quad \mathrm{color}(z)\ne\mathrm{color}(u),\quad
 (u\co z\co v\ \lor\ v\co z\co u).
\]
Project $z$, complement, and intersect with equal-color endpoints in $D$.
This computes $E$. These are standard automatic-structure
operations~\cite{Blum99,BG00}. Next compute the regular language
\begin{equation}\label{eq:rep}
 \Rep(u)\iff u\in D\ \land\
 \neg\exists v\in D\;(E(u,v)\land
            \rev(v)<_{\mathrm{shortlex}}\rev(u)).
\end{equation}
Shortlex is a well-order on a finite alphabet, so it selects exactly one
representative per component, including components without a least colex
word. Remove padding for the ordinary word encoding. A finite automaton
decides finiteness of $\Rep$: its reachable, coaccessible subgraph contains
a genuine-letter cycle exactly when the language is infinite. If finite,
enumerate its words in the useful acyclic graph. Thus the number of convex
components is decidable without a guessed bound. Fixing a representative
in $E$ also gives an effective regular language for its component.

The relation $E$ is right invariant. If $uEv$, their $C$-colors give the same
defined letters and successor colors. For $u\co v$, any word $z\in D$
between $ua$ and $va$ ends in $a$, because that suffix cylinder is convex.
Write $z=wa$. Prefix closure puts $w$ in $D$, and suffix cancellation gives
$u\co w\co v$. Hence $w$ has the endpoint color and $z$ the common
successor color. Therefore $ua\,E\,va$. Equality and the reverse endpoint
order give the other cases; partial definedness also agrees.

If $\Rep$ is finite, build a transducer $C_c$ with one state per component,
initial component $\{\eps\}$, and the transition and terminal outputs of
its $C$-color. At representative $u$, the target on $a$ is the unique
component whose representative $r$ satisfies $E(ua,r)$. Right invariance
makes the choice independent of $u$. Induction shows that every $w\in D$
reaches its own component with exactly $A_C(w)$, and the same accepted
domain and terminal output as $C$. Every component is accessible and
coaccessible. Thus $C_c$ is finite, trimmed, equivalent to $f$, and its
reaching languages are exactly the separated convex components.

\subsection{Why no other realization can evade the tests}

Suppose a finite Wheeler $W$ computes $f$. Its reaching sets $J_p$ are
convex in $D$. Each $J_p\setminus\{\eps\}$ is canonically monochromatic by
\eqref{eq:otherH}, and removing epsilon preserves convexity because epsilon
is least. These sets, with singleton epsilon, give a finite convex
refinement of the canonical coloring. Therefore $E$ has finitely many
classes and $\Rep$ is finite.

Moreover, common right append of $\lambda'_p$ preserves each orientation
on each monotone piece of $A_W|_{J_p\setminus\{\eps\}}$. Equation
\eqref{eq:otherH} transfers those pieces to $H$. Finitely many states and
pieces thus give a finite convex monotone partition of $H$ on
$D\setminus\{\eps\}$. Add singleton epsilon with output $A_C(\eps)=\eps$.
By~\eqref{eq:H} this is a partition for $A_C$ on $D$. Intersect it with
each convex $E$-class. The resulting finitely many convex monotone pieces
prove the output condition for $C_c$, whose input order already holds.
Consequently \emph{every} Wheeler realization forces this particular
canonical convex split to pass membership.

The complete decision algorithm is now: construct $C$; compute $E$ and
$\Rep$; reject if $\Rep$ is infinite; otherwise construct $C_c$ and apply
Theorem~\ref{thm:member}, returning $C_c$ exactly on acceptance. All steps
terminate. An infinite representative language rules out every Wheeler
realization; a failed membership test also rules out every such realization
by the preceding necessity argument. A positive answer supplies an
equivalent Wheeler machine. This proves Theorem~\ref{thm:realize}.

\section{Verification, scope, and AI methodology}

The proofs establish the claims for all promised finite inputs. Deterministic
Python prototypes additionally implement regular-range comparison, seed
propagation, SCC tests, normalization and residual merging, synchronous
projection, regular representatives, finiteness, and the split construction.
The membership controls covered all 544 accessible input-ordered binary
partial machines through two states with outputs in $\{\eps,0,1\}$, and
600 seeded machines with three to five states. All states were final.
Direct sorted finite paths and a segmentation dynamic program checked the
reported implications and piece bounds, through lengths seven and eight.
Realizability controls included 13 boundary examples and 20 complete finite
domains, checking 1,475 function values, 1,318 LCP identities, 8,171 relation
comparisons, and 1,315 component-successor tests. Complete finite domains
were compared against their entire direct convex decompositions; infinite
controls used exact regular-language finiteness with supplementary literal
replays through length seven.

Checks included erasure, partial terminal domains, the epsilon-only domain,
redundant unary parity, nonempty initial LCP with an initial cycle, and
finite canonical components whose output still fails membership. A separate
model context reconstructed the arguments and ran additional finite checks.
These activities are internal computational and mathematical checks, not
external peer review or formal verification. No absence of a sampled
counterexample is used as an infinite assertion. The accompanying local
code and reports specify the finite coverage and retain failed approaches.

The membership exponent is not optimized, the certificate bound is only
computable, and no practical indexing improvement, minimal realization, or
realizability complexity bound is asserted. The results concern the stated
piecewise definition with its literal reaching-language convention.
The literature comparison is bounded by inspected primary sources; it is
not a universal priority claim.

\paragraph{AI methodology.}
OpenAI Codex was used substantially throughout the project: it proposed and
developed proof arguments, searched primary literature, generated and executed
verification code, performed separate-context mathematical checks, and drafted
and revised the exposition. This materially affected all technical sections,
the computational controls, and the references. The checks do not constitute
formal proofs or independent human peer review. No claim of an additional
human verification is made here; the named author retains responsibility for
the content and any future submission.

\clearpage
\bibliographystyle{ACM-Reference-Format}
\bibliography{references}
\end{document}
